\chapter{The example in numbers}\label{app:numbers}

The earlier chapters work in symbols. This chapter collects the values of the paper's example in one place.
Use it only when you read the three figures, or when you want to see the size of an effect.
Every value comes from the paper. Money is in units of the low value \(\ell\).

\section{The primitives}

\begin{center}
\begin{tabular}{@{}llc@{}}
\toprule
Object & Symbol & Value\\
\midrule
High-class value & \(h\) & \(\reg{prim_h}\)\\
Low-class value & \(\ell\) & \(\reg{prim_ell}\)\\
Reserve & \(p\) & \(\reg{prim_p}\)\\
Prior of class \(H\) & \(\pi\) & \(\reg{prim_pi}\)\\
Probability of a zero cost & \(\rho\) & \(\reg{prim_rho}\)\\
High preparation cost & \(c_H\) & \(\reg{prim_cH}\)\\
Trading cost & \(k\) & \(\reg{prim_k}\)\\
Weak, strong and very strong incumbent & \(r_0,\ r_1,\ r_2\) & \(\reg{prim_r0},\ \reg{prim_r1},\ \reg{prim_r2}\)\\
\midrule
Largest likelihood ratio of one order & \(1+u\) & \(\eu\)\\
Liquidity probability & \(\beta\) & \(\reg{beta_B}\)\\
Presence of the informed investor & \(\lambda\) & \(\reg{lambda_B}\)\\
Liquidity order in one direction & \(b\) & \(\reg{b_B}\)\\
Top and bottom posterior at \(\pi=1/2\) & \(M(1/2),\ m(1/2)\) & \(\reg{M_half_B},\ \reg{m_half_B}\)\\
\bottomrule
\end{tabular}
\end{center}

\section{Payoffs and outcomes at the three strengths}

\begin{center}
\begin{tabular}{@{}lccc@{}}
\toprule
 & \(r_0\) & \(r_1\) & \(r_2\)\\
\midrule
Target-payoff spread \(\DT\) & \(\reg{Delta_r0}\) & \(\reg{Delta_r1}\) & \(\reg{Delta_r2}\)\\
Gross profit at the prior \(B_r(1/2)\) & \(\reg{Bprior_r0}\) & \(\reg{Bprior_r1_B}\) & \(\reg{Bprior_r2}\)\\
Preparation threshold \(\tau(r)\) & \(\reg{tau_r0}\) & \(\reg{tau_r1}\) & \(\reg{tau_r2}\)\\
Entry \(\Ent\) & \(\reg{entry_r0}\) & \(\reg{entry_r1_B}\) & \(\reg{entry_r2}\)\\
High-class ownership \(\Own_H\) & \(\reg{own_r0}\) & \(\reg{own_r1_B}\) & \(\reg{own_r2_B}\)\\
\bottomrule
\end{tabular}
\end{center}

From \(r_0\) to \(r_1\), gross profit at the prior falls while the spread rises. Entry rises and then returns to the floor \(\rho\) at \(r_2\).

\section{The band of Corollary 1}

At the prior \(1/2\), the conditions (H1) to (H3) hold exactly when \(u\) lies in \((\reg{u_lo},\reg{u_hi})\).
In terms of the top posterior, the band is \(\tau(r_1)=\reg{M_lo}<M(1/2)<\reg{M_hi}=\tau(r_2)\).
The example's slot has \(M(1/2)=\reg{M_half_B}\), inside the band.
With \(1+u=\eu\), the prior must lie between about \(\reg{prior_lo_B}\) and \(\reg{prior_hi_B}\).
Class \(H\) has a decisive advantage because \(c_H\) exceeds \((h^2-p^2)/(2h)=\reg{cH_DA_floor_B}\).

\section{Marked strengths in the figures}

\begin{center}
\begin{tabular}{@{}llc@{}}
\toprule
Strength & Where it lies & Value\\
\midrule
\multicolumn{3}{@{}l}{\emph{Coexistence (Chapter~\ref{ch:coexistence}, Figure of entry along the path)}}\\
\(r_A\) & start of coexistence & \(\reg{rA}\)\\
\(r_F\) & class \(H\) buys for sure from here & \(\reg{rF_B}\)\\
\(r_N\) & end of the quiet market & \(\reg{rN}\)\\
\(r_J\) & end of mixed orders at the entry floor & \(\reg{rJ}\)\\
\(r_S\) & class \(L\) sells for sure from here & \(\reg{rS_B}\)\\
\(r_C\) & entry ceiling, where \(\tau(r)=M(1/2)\) & \(\reg{rC_B}\)\\
\midrule
\multicolumn{3}{@{}l}{\emph{Research (Chapter~\ref{ch:research}), with \(a=\reg{res_a}\) and cap \(\lb\approx\reg{res_cap}\)}}\\
\(r_R\) & recruiting research starts & \(\reg{res_rR}\)\\
\(r_{\lb}\) & research reaches the cap & \(\reg{res_rcap}\)\\
\(r_N\) & zero research ends & \(\reg{res_rN}\)\\
\(r_Q\) & floor-entry research below the cap ends & \(\reg{res_rQ}\)\\
\(r_C\) & recruiting research ends & \(\reg{res_rC}\)\\
\bottomrule
\end{tabular}
\end{center}

Entry with recruiting research rises from \(\reg{res_entry_rR}\) at \(r_R\) to \(\reg{res_entry_cap}\) at \(r_{\lb}\).
